\begin{pastexam}[label={px:ntu-114-bime-ds-06}]{臺大生機所 114 學年度 第 6 題}
（問答；10 分）

\textbf{(10\%)} Consider a sorted circular doubly-linked list where the head element points to the
smallest element in the list.\\
\textbf{a.} What is the complexity of determining whether a given element exists in the list?\\
\textbf{b.} What is the complexity of deleting a given element e in the list (not including the
cost of finding it)?
\end{pastexam}

\begin{pastsoln}
設串列有 $n$ 個元素，節點有 $\mathit{prev}$、$\mathit{next}$，$\mathit{head}$ 指向最小的元素；
環狀，所以 $\mathit{head}.\mathit{prev}$ 是最大的元素（資結書第~5.2~節）。

\textbf{a.} $O(n)$（最壞 $\Theta(n)$）。 串列沒有隨機存取，要找到第 $i$ 個節點只能從
$\mathit{head}$ 一步步走過去，所以排好序也不能二分搜尋（光是走到中間就要 $n/2$ 步）。 排序和環狀
可以幫一點忙：先看 $\mathit{head}$（最小）和 $\mathit{head}.\mathit{prev}$（最大），超出範圍就直接
回答不在，$O(1)$；從頭往後走時，碰到比目標大的鍵值就可以停，不必走完。 雙向也讓我們能從比較近的一
端走：目標靠近最大值時從 $\mathit{head}.\mathit{prev}$ 往回走。 但這些都只省常數倍；目標在中間
時，不論從哪一端都要走大約 $n/2$ 個節點，最壞仍是 $\Theta(n)$。

\textbf{b.} $O(1)$。 已經拿到 $e$ 的指標，雙向串列可以直接從 $e$ 找到前一個和後一個節點，不必像
單向串列那樣從頭找前一個節點：
\begin{pseudo}{刪除環狀雙向串列裡的節點 $e$}
\Procedure{Delete}{$L, e$}
    \If{$e.\mathit{next} = e$}
        \State $L.\mathit{head} \gets \NIL$ \Comment{只有 $e$ 一個節點}
    \Else
        \State $e.\mathit{prev}.\mathit{next} \gets e.\mathit{next}$
        \State $e.\mathit{next}.\mathit{prev} \gets e.\mathit{prev}$
        \If{$L.\mathit{head} = e$}
            \State $L.\mathit{head} \gets e.\mathit{next}$ \Comment{新的最小值}
        \EndIf
    \EndIf
\EndProcedure
\end{pseudo}
固定幾個指標操作，和 $n$ 無關，所以是 $O(1)$。 刪掉一個元素不會破壞排序；刪掉的若是最小值，
$\mathit{head}$ 改指向下一個，它就是新的最小值。
\end{pastsoln}
